% Lab: Specific Heat Capacity of Aluminum & Leakiness Test

% Purpose:
% To find the specific heat capacity of aluminum by experiment & compare the experimental value to reference material
% Materials:

% Equipment Chemicals

% ● thermometer ● distilled water (at room temperature) for the calorimeter
% ● three Styrofoam cups + lid ● aluminum pellets
% ● one hot plate ● tap water (any temperature) for the hot water bath
% ● test tube with rubber stopper
% ● one beaker (250 mL)
% ● one graduated cylinder (100 mL)
% ● mass scale
% ● time piece (i.e., watch, timer, phone)
% Procedure:
% 1. Fill a beaker about two thirds full with tap water. Start heating it on the hot plate.
% 2. Obtain about 15 g of aluminum. Record the exact mass.
% 3. Place the aluminum in a test tube. Place the rubber stopper LOOSELY on top. Place the test tube & contents into the
% hot water bath. Make sure the metal is sitting below the water’s surface. Heat until the water boils. Continue heating
% for at least another 10 minutes to ensure the metal attains the temperature of the boiling water.
% 4. While the water boils, take 3 Styrofoam cups. Clean & dry them if necessary. Place one in the other. This becomes
% the calorimeter.
% 5. Weigh the empty calorimeter.
% 6. Using the graduated cylinder, place about 30 mL of distilled water into the calorimeter & weigh it again.
% 7. Obtain the temperature of the water in the calorimeter to the nearest 0.1°C. The thermometer bulb should be
% completely under the water. Hold the thermometer at the very top to minimize thermal energy transfer from your hand
% to the thermometer’s contents. Do not leave the thermometer in the calorimeter! It can tip over!
% 8. When ready, take the test tube out of the boiling water, remove the stopper & pour the aluminum into the water of the
% calorimeter. The moment the pellets touch the water, start the timer on a time piece. Be careful that any water
% adhering to the test tube does not fall into the calorimeter. Replace the calorimeter cover & agitate the water the best
% you can GENTLY with the thermometer (it is better to use a stir-rod to stir, but if you are careful, the thermometer will
% suffice). Record to the nearest 0.1°C the maximum temperature reached by the water. Record how long it took for the
% highest temperature to be achieved. After that, wait 3 min (180.0 s) & then record the temperature again.
% 9. Repeat steps 5 – 8. Be sure to dry your metal before reusing it; this can be done by heating the metal briefly in the test
% tube in boiling water & then pouring the metal onto a paper towel to drain. (Time permitting, repeat the experiment one
% more time, for a total of 3 trials). For trials 2 & 3, do not perform the 3 min wait.
% Calculations Regarding Time:
% Calculate the rate at which the temperature drops during the 180.0 s (3 minutes) after the maximum temperature is reached
% using the equation:

% rate =
% final temp − temp after 3 min
% 180 seconds

% This value reflects the rate at which thermal energy leaks out of the calorimeter (in °C/s). This is how “leaky” the calorimeter is at
% the end. Assuming this thermal energy rate loss is constant throughout the process, multiply this rate by the number of seconds
% between the initial temperature & the final temperature. If necessary, add this time correction to the measured maximum
% temperature to obtain a corrected maximum temperature.
% If the corrected temperature is not significantly different from the measured one

% (i.e.,within 5% of the measured final temperature, ௖௢௥௥௘௖௧௘ௗି௠௘௔௦௨௥௘ௗ
% ௠௘௔௦௨௥௘ௗ

% < 5%)

% then ignore this correction for the other parts of the lab & use only the measured final temperatures for all parts of the lab
% analysis. If the difference between the corrected temperature & the measured temperature is more than 5%, use the corrected
% temperature for all parts of the experimental analysis. Assume the corrected temperature is the correct value.

% Observations/Results:
% Design your own observations table(s). Be sure to leave enough space to record relevant data from each trial. Any additional
% observations are welcome.
% Analysis & Calculations:
% From the results, calculate the average specific heat capacity for aluminum & determine the % error. Make sure the
% calculations are clear & easy-to-follow. Use headings if necessary to emphasize clarity. Use proper problem-solving format.

% Prelab Instructions:
% Have the following completed before the lab date:
% ● Date, Title & Purpose
% ● Introduction
% ● Observations Table(s) drawn & ready to go

% Postlab Instructions:
% In addition to the prelab, the completed lab report will have the
% following:
% ● Lab partner’s name(s)
% ● Completed observation(s) table(s)
% ● Analysis
% ● Discussion Question(s)


\documentclass{article}
\usepackage[margin=1in]{geometry}
\usepackage{tikz}
\usepackage{amsmath}
\usepackage{amssymb}
\usepackage{multicol}
\usepackage{import}

\import{../}{labflowchart.sty}
\renewcommand{\thesection}{\Roman{section}} 
\renewcommand{\thesubsection}{\thesection.\Roman{subsection}}
\begin{document}
\title{Lab: Specific Heat Capacity of Aluminum \& Leakiness Test}
\author{Aaron Huang}
\date{March 25, 2026}
\maketitle
\section{Purpose}
To find the specific heat capacity of aluminum by experiment \& compare the experimental value to reference material
\section{Introduction}
% INTRODUCTION
% • write 3 short paragraphs using complete sentences
% • it includes:
% • paragraph #1: describes/explains the concept/theory of the lab (2 – 5 sentences)
% • paragraph #2: an overview of the lab procedure (2 – 5 sentences, no need to indicate specific amounts)
% • paragraph #3: special safety considerations regarding the use of specific chemicals &/or equipment beyond the
% necessity of wearing goggles (1 – 2 sentences) (sometimes there are no special safety considerations)
This lab was designed to determine the specific heat capacity of aluminum by experiment using calorimetry. The specific heat capacity of a substance is the amount of energy required to raise the temperature of one gram of the substance by one degree Celsius. This experiment, used a calorimeter, a device that minimizes heat and mass transfer between the system and its surroundings, to measure the heat transfer between aluminum and water. Additionally, the experiment included a leakiness test to determine the rate at which thermal energy leaks out of the calorimeter, which is important for accurate measurements.

The procedure involved heating aluminum pellets in a hot water bath until they reached the temperature of boiling water, 100°C, then the aluminum was transferred to a calorimeter containing water with known mass and initial temperature. The maximum temperature reached by the water was recorded, as well as the time taken to reach that temperature. Finally, after waiting for 3 minutes, the temperature of the water once again was recorded to determine the rate of heat loss from the calorimeter. 

% Safty regarding hot plate
The hot plate, boing water, and aluminum pellets were handled with care to avoid burns and the rubber stopper was put on loosely to prevent pressure buildup. Additionally, the hot plate was not moved while in use and unplugged after use to prevent accidents.

\section{Observations and Results}

% This command increases the vertical padding in table cells
\renewcommand{\arraystretch}{1.6} 

\begin{table}[h]
\centering
\caption{Experimental Data for Specific Heat of Aluminum}
\begin{tabular}{|p{6cm}|c|c|c|c|}
\hline
\textbf{Measurement} & \textbf{Unit} & \textbf{Trial 1} & \textbf{Trial 2} & \textbf{Trial 3} \\ \hline
Mass of empty calorimeter & g & & & \\ \hline
Mass of calorimeter + water & g & & & \\ \hline
Mass of aluminum pellets ($m_{Al}$) & g & & & \\ \hline
Initial temp of water ($T_{i,w}$) & $^\circ$C & & & \\ \hline
Initial temp of Al ($T_{i,Al}$) & $^\circ$C & & & \\ \hline
Maximum temp reached ($T_{max}$) & $^\circ$C & & & \\ \hline
Time to reach $T_{max}$ ($\Delta t$) & s & & --- & --- \\ \hline
Temp after 3-minute wait ($T_{3min}$) & $^\circ$C & & --- & --- \\ \hline
\end{tabular}
\end{table}

\renewcommand{\arraystretch}{1.0} % Resets row height for the rest of the doc


\end{document}